% 概率论与数理统计例题回顾 —— 正文排版集中在 0.样式.tex
\documentclass[11pt,a4paper,oneside,openany]{book}
\usepackage[UTF8,fontset=windows,scheme=plain]{ctex}
\usepackage{fontspec}
\setmainfont{Times New Roman}
\setsansfont{Microsoft YaHei}
\newcommand{\YHBbooktrue}{}
\newcommand{\YHBcourse}{概率论与数理统计}
\newcommand{\YHBversion}{v1.1}
\newcommand{\YHBquote}{一以贯之的努力，不得懈怠的人生。\par\vspace{2mm}\normalfont —— 《天才基本法》}
\newcommand{\YHBcoverSize}{34pt}
\newcommand{\YHBcoverLead}{42pt}
\newcommand{\YHBcoversub}{例题回顾}
\newcommand{\YHBdoclabel}{例题回顾}
\usepackage{xcolor}
\definecolor{yhbAccent}{HTML}{3A2E5C}
\definecolor{yhbAccent2}{HTML}{6B5B8C}
% ============================================================
% 0.样式.tex —— 课程笔记统一正文排版
% 用法：主文件先定义课程信息与主色，再 % ============================================================
% 0.样式.tex —— 课程笔记统一正文排版
% 用法：主文件先定义课程信息与主色，再 % ============================================================
% 0.样式.tex —— 课程笔记统一正文排版
% 用法：主文件先定义课程信息与主色，再 \input{0.样式.tex}
%   \newcommand{\YHBcourse}{微积分A}      % 课程名（页眉与书签用）
%   \newcommand{\YHBversion}{v1.2}        % 版本号（封面用）
%   \newcommand{\YHBquote}{……}            % 封面题句
%   \definecolor{yhbAccent}{HTML}{24465C} % 主色（深）
%   \definecolor{yhbAccent2}{HTML}{347A80}% 辅色（浅）
%   \YHBbooktrue / \YHBbookfalse          % 是否 book 文档类（有 \chapter）
% 章节文件只写内容，不要重复 \usepackage、不要改字号。
% ============================================================

% ---------- 页面与字体 ----------
\usepackage{geometry}
\geometry{a4paper,left=24mm,right=24mm,top=23mm,bottom=24mm,headheight=15pt,headsep=7mm,footskip=13mm}
\usepackage{microtype}
\usepackage{amsmath,amssymb,amsfonts,mathtools}
\usepackage{esint}
\usepackage{graphicx}
\usepackage{adjustbox}
\usepackage{array,booktabs,multirow,makecell,tabularx,longtable}
\usepackage{xcolor,colortbl}
\usepackage{enumitem}
\usepackage{needspace,etoolbox}
\usepackage{titlesec,titletoc}
\usepackage{fancyhdr}
\usepackage{caption}
\usepackage[most]{tcolorbox}

% ---------- 颜色别名（章节里用的旧名继续有效） ----------
\colorlet{mainblue}{yhbAccent}
\colorlet{secondcyan}{yhbAccent2}
\colorlet{thirdteal}{yhbAccent2}
\colorlet{highlight}{yhbAccent!6!white}
\colorlet{tablehead}{yhbAccent!12!white}
\definecolor{textgray}{HTML}{4A5057}

% ---------- 课程标签 ----------
\providecommand{\YHBdoclabel}{课程笔记}
% ---------- 正文节奏 ----------
\linespread{1.3}
\setlength{\parindent}{0pt}
\setlength{\parskip}{6pt plus 1.5pt minus 1.5pt}
\setlength{\emergencystretch}{2.5em}
\raggedbottom
\widowpenalty=10000
\clubpenalty=10000
\setlength{\tabcolsep}{7pt}
\renewcommand{\arraystretch}{1.3}
\allowdisplaybreaks[2]
\AtBeginDocument{%
  \setlength{\abovedisplayskip}{10pt plus 2pt minus 2pt}%
  \setlength{\belowdisplayskip}{10pt plus 2pt minus 2pt}%
  \setlength{\abovedisplayshortskip}{6pt}%
  \setlength{\belowdisplayshortskip}{6pt}%
}
\captionsetup{font=small,labelfont=bf,skip=6pt}

% ---------- 列表 ----------
\setlist{leftmargin=2em,labelsep=0.6em,itemsep=4pt,parsep=2pt,topsep=6pt,partopsep=0pt}
\setlist[itemize,2]{label=\textbullet,leftmargin=1.6em}

% ---------- 标题层级 ----------
\ifdefined\YHBbooktrue
  \titleformat{\chapter}[display]
    {\sffamily\bfseries\color{yhbAccent}}
    {\normalsize\sffamily 第\,\thechapter\,章}
    {6pt}
    {\LARGE}
    [{\vspace{5pt}{\color{yhbAccent!55}\titlerule[1pt]}}]
  \titlespacing*{\chapter}{0pt}{14pt}{24pt}
  \titleformat{\section}{\Large\sffamily\bfseries\color{yhbAccent}}{\thesection}{0.7em}{}
  \titlespacing*{\section}{0pt}{18pt plus 3pt minus 2pt}{8pt}
\else
  \titleformat{\section}{\Large\sffamily\bfseries\color{yhbAccent}}{\thesection}{0.7em}{}[{\vspace{3pt}{\color{yhbAccent!45}\titlerule[0.7pt]}}]
  \titlespacing*{\section}{0pt}{18pt plus 3pt minus 2pt}{9pt}
\fi
\titleformat{\subsection}{\large\sffamily\bfseries\color{yhbAccent!85!black}}{\thesubsection}{0.7em}{}
\titleformat{\subsubsection}{\normalsize\bfseries\color{black!85}}{\thesubsubsection}{0.7em}{}
\titlespacing*{\subsection}{0pt}{12pt plus 2pt minus 1pt}{5pt}
\titlespacing*{\subsubsection}{0pt}{9pt plus 2pt minus 1pt}{4pt}
\pretocmd{\section}{\Needspace{6\baselineskip}}{}{}
\pretocmd{\subsection}{\Needspace{5\baselineskip}}{}{}
\pretocmd{\subsubsection}{\Needspace{4\baselineskip}}{}{}

\ifdefined\YHBbooktrue
  \setcounter{secnumdepth}{2}
  \setcounter{tocdepth}{1}
\else
  \setcounter{secnumdepth}{3}
  \setcounter{tocdepth}{2}
\fi
\renewcommand{\contentsname}{目录}
% 目录编号颜色由 titletoc 的 titlecontents 控制，不再重定义 \numberline

% ---------- 目录样式 ----------
\ifdefined\YHBbooktrue
  \titlecontents{chapter}[0pt]
    {\addvspace{6pt}\sffamily\bfseries\color{yhbAccent}}
    {\contentslabel{2.2em}}{}{\titlerule*[0.6em]{.}\contentspage}
  \titlecontents{section}[2.2em]
    {\addvspace{2pt}}
    {\contentslabel{2.2em}}{}{\titlerule*[0.6em]{.}\contentspage}
\else
  \titlecontents{section}[0pt]
    {\addvspace{5pt}\sffamily\bfseries\color{yhbAccent}}
    {\contentslabel{2.2em}}{}{\titlerule*[0.6em]{.}\contentspage}
  \titlecontents{subsection}[2.2em]
    {\addvspace{2pt}}
    {\contentslabel{2.4em}}{}{\titlerule*[0.6em]{.}\contentspage}
\fi

% ---------- 页眉页脚 ----------
\pagestyle{fancy}
\fancyhf{}
\fancyhead[L]{\footnotesize\sffamily\color{textgray}\YHBcourse\ \YHBdoclabel}
\ifdefined\YHBbooktrue
  \fancyhead[R]{\footnotesize\sffamily\color{textgray}\nouppercase{\leftmark}}
  \renewcommand{\chaptermark}[1]{\markboth{#1}{}}
\else
  \fancyhead[R]{\footnotesize\sffamily\color{textgray}\nouppercase{\rightmark}}
  \renewcommand{\sectionmark}[1]{\markright{#1}}
\fi
\fancyfoot[C]{\footnotesize\color{textgray}\thepage}
\renewcommand{\headrulewidth}{0.35pt}
\renewcommand{\headrule}{{\color{yhbAccent!45}\hrule width\headwidth height\headrulewidth}}
\fancypagestyle{plain}{%
  \fancyhf{}%
  \fancyfoot[C]{\footnotesize\color{textgray}\thepage}%
  \renewcommand{\headrulewidth}{0pt}%
}

% ---------- 公式与考点框 ----------
\newcommand{\keyformula}[1]{%
  \begin{tcolorbox}[enhanced,breakable,colback=highlight,colframe=yhbAccent!30!white,boxrule=0.5pt,arc=2.5pt,left=11pt,right=11pt,top=9pt,bottom=9pt,boxsep=0pt,before skip=11pt,after skip=11pt,halign=center,width=\linewidth]
  \adjustbox{max width=\linewidth-2em}{$\displaystyle #1$}
  \end{tcolorbox}}
\newcommand{\corepoint}[1]{%
  \begin{tcolorbox}[enhanced,breakable,colback=highlight,colframe=yhbAccent2,boxrule=0pt,borderline west={2.4pt}{0pt}{yhbAccent2},arc=0pt,left=11pt,right=11pt,top=8pt,bottom=8pt,before skip=11pt,after skip=11pt]
  \textbf{\color{yhbAccent}核心考点}\quad #1
  \end{tcolorbox}}
\newcommand{\symb}[2]{\textbf{#1}：#2\par}

% ---------- 超链接（务必最后加载） ----------
\usepackage[unicode=true,colorlinks=true,linkcolor=yhbAccent,citecolor=yhbAccent,urlcolor=yhbAccent!85!black,pdfborder={0 0 0},bookmarksnumbered=true,pdfstartview=FitH]{hyperref}
\usepackage{hypcap}
\hypersetup{pdftitle={\YHBcourse 课程笔记},pdfauthor={浮萍},pdfsubject={课程学习与复习笔记},pdfcreator={XeLaTeX}}

%   \newcommand{\YHBcourse}{微积分A}      % 课程名（页眉与书签用）
%   \newcommand{\YHBversion}{v1.2}        % 版本号（封面用）
%   \newcommand{\YHBquote}{……}            % 封面题句
%   \definecolor{yhbAccent}{HTML}{24465C} % 主色（深）
%   \definecolor{yhbAccent2}{HTML}{347A80}% 辅色（浅）
%   \YHBbooktrue / \YHBbookfalse          % 是否 book 文档类（有 \chapter）
% 章节文件只写内容，不要重复 \usepackage、不要改字号。
% ============================================================

% ---------- 页面与字体 ----------
\usepackage{geometry}
\geometry{a4paper,left=24mm,right=24mm,top=23mm,bottom=24mm,headheight=15pt,headsep=7mm,footskip=13mm}
\usepackage{microtype}
\usepackage{amsmath,amssymb,amsfonts,mathtools}
\usepackage{esint}
\usepackage{graphicx}
\usepackage{adjustbox}
\usepackage{array,booktabs,multirow,makecell,tabularx,longtable}
\usepackage{xcolor,colortbl}
\usepackage{enumitem}
\usepackage{needspace,etoolbox}
\usepackage{titlesec,titletoc}
\usepackage{fancyhdr}
\usepackage{caption}
\usepackage[most]{tcolorbox}

% ---------- 颜色别名（章节里用的旧名继续有效） ----------
\colorlet{mainblue}{yhbAccent}
\colorlet{secondcyan}{yhbAccent2}
\colorlet{thirdteal}{yhbAccent2}
\colorlet{highlight}{yhbAccent!6!white}
\colorlet{tablehead}{yhbAccent!12!white}
\definecolor{textgray}{HTML}{4A5057}

% ---------- 课程标签 ----------
\providecommand{\YHBdoclabel}{课程笔记}
% ---------- 正文节奏 ----------
\linespread{1.3}
\setlength{\parindent}{0pt}
\setlength{\parskip}{6pt plus 1.5pt minus 1.5pt}
\setlength{\emergencystretch}{2.5em}
\raggedbottom
\widowpenalty=10000
\clubpenalty=10000
\setlength{\tabcolsep}{7pt}
\renewcommand{\arraystretch}{1.3}
\allowdisplaybreaks[2]
\AtBeginDocument{%
  \setlength{\abovedisplayskip}{10pt plus 2pt minus 2pt}%
  \setlength{\belowdisplayskip}{10pt plus 2pt minus 2pt}%
  \setlength{\abovedisplayshortskip}{6pt}%
  \setlength{\belowdisplayshortskip}{6pt}%
}
\captionsetup{font=small,labelfont=bf,skip=6pt}

% ---------- 列表 ----------
\setlist{leftmargin=2em,labelsep=0.6em,itemsep=4pt,parsep=2pt,topsep=6pt,partopsep=0pt}
\setlist[itemize,2]{label=\textbullet,leftmargin=1.6em}

% ---------- 标题层级 ----------
\ifdefined\YHBbooktrue
  \titleformat{\chapter}[display]
    {\sffamily\bfseries\color{yhbAccent}}
    {\normalsize\sffamily 第\,\thechapter\,章}
    {6pt}
    {\LARGE}
    [{\vspace{5pt}{\color{yhbAccent!55}\titlerule[1pt]}}]
  \titlespacing*{\chapter}{0pt}{14pt}{24pt}
  \titleformat{\section}{\Large\sffamily\bfseries\color{yhbAccent}}{\thesection}{0.7em}{}
  \titlespacing*{\section}{0pt}{18pt plus 3pt minus 2pt}{8pt}
\else
  \titleformat{\section}{\Large\sffamily\bfseries\color{yhbAccent}}{\thesection}{0.7em}{}[{\vspace{3pt}{\color{yhbAccent!45}\titlerule[0.7pt]}}]
  \titlespacing*{\section}{0pt}{18pt plus 3pt minus 2pt}{9pt}
\fi
\titleformat{\subsection}{\large\sffamily\bfseries\color{yhbAccent!85!black}}{\thesubsection}{0.7em}{}
\titleformat{\subsubsection}{\normalsize\bfseries\color{black!85}}{\thesubsubsection}{0.7em}{}
\titlespacing*{\subsection}{0pt}{12pt plus 2pt minus 1pt}{5pt}
\titlespacing*{\subsubsection}{0pt}{9pt plus 2pt minus 1pt}{4pt}
\pretocmd{\section}{\Needspace{6\baselineskip}}{}{}
\pretocmd{\subsection}{\Needspace{5\baselineskip}}{}{}
\pretocmd{\subsubsection}{\Needspace{4\baselineskip}}{}{}

\ifdefined\YHBbooktrue
  \setcounter{secnumdepth}{2}
  \setcounter{tocdepth}{1}
\else
  \setcounter{secnumdepth}{3}
  \setcounter{tocdepth}{2}
\fi
\renewcommand{\contentsname}{目录}
% 目录编号颜色由 titletoc 的 titlecontents 控制，不再重定义 \numberline

% ---------- 目录样式 ----------
\ifdefined\YHBbooktrue
  \titlecontents{chapter}[0pt]
    {\addvspace{6pt}\sffamily\bfseries\color{yhbAccent}}
    {\contentslabel{2.2em}}{}{\titlerule*[0.6em]{.}\contentspage}
  \titlecontents{section}[2.2em]
    {\addvspace{2pt}}
    {\contentslabel{2.2em}}{}{\titlerule*[0.6em]{.}\contentspage}
\else
  \titlecontents{section}[0pt]
    {\addvspace{5pt}\sffamily\bfseries\color{yhbAccent}}
    {\contentslabel{2.2em}}{}{\titlerule*[0.6em]{.}\contentspage}
  \titlecontents{subsection}[2.2em]
    {\addvspace{2pt}}
    {\contentslabel{2.4em}}{}{\titlerule*[0.6em]{.}\contentspage}
\fi

% ---------- 页眉页脚 ----------
\pagestyle{fancy}
\fancyhf{}
\fancyhead[L]{\footnotesize\sffamily\color{textgray}\YHBcourse\ \YHBdoclabel}
\ifdefined\YHBbooktrue
  \fancyhead[R]{\footnotesize\sffamily\color{textgray}\nouppercase{\leftmark}}
  \renewcommand{\chaptermark}[1]{\markboth{#1}{}}
\else
  \fancyhead[R]{\footnotesize\sffamily\color{textgray}\nouppercase{\rightmark}}
  \renewcommand{\sectionmark}[1]{\markright{#1}}
\fi
\fancyfoot[C]{\footnotesize\color{textgray}\thepage}
\renewcommand{\headrulewidth}{0.35pt}
\renewcommand{\headrule}{{\color{yhbAccent!45}\hrule width\headwidth height\headrulewidth}}
\fancypagestyle{plain}{%
  \fancyhf{}%
  \fancyfoot[C]{\footnotesize\color{textgray}\thepage}%
  \renewcommand{\headrulewidth}{0pt}%
}

% ---------- 公式与考点框 ----------
\newcommand{\keyformula}[1]{%
  \begin{tcolorbox}[enhanced,breakable,colback=highlight,colframe=yhbAccent!30!white,boxrule=0.5pt,arc=2.5pt,left=11pt,right=11pt,top=9pt,bottom=9pt,boxsep=0pt,before skip=11pt,after skip=11pt,halign=center,width=\linewidth]
  \adjustbox{max width=\linewidth-2em}{$\displaystyle #1$}
  \end{tcolorbox}}
\newcommand{\corepoint}[1]{%
  \begin{tcolorbox}[enhanced,breakable,colback=highlight,colframe=yhbAccent2,boxrule=0pt,borderline west={2.4pt}{0pt}{yhbAccent2},arc=0pt,left=11pt,right=11pt,top=8pt,bottom=8pt,before skip=11pt,after skip=11pt]
  \textbf{\color{yhbAccent}核心考点}\quad #1
  \end{tcolorbox}}
\newcommand{\symb}[2]{\textbf{#1}：#2\par}

% ---------- 超链接（务必最后加载） ----------
\usepackage[unicode=true,colorlinks=true,linkcolor=yhbAccent,citecolor=yhbAccent,urlcolor=yhbAccent!85!black,pdfborder={0 0 0},bookmarksnumbered=true,pdfstartview=FitH]{hyperref}
\usepackage{hypcap}
\hypersetup{pdftitle={\YHBcourse 课程笔记},pdfauthor={浮萍},pdfsubject={课程学习与复习笔记},pdfcreator={XeLaTeX}}

%   \newcommand{\YHBcourse}{微积分A}      % 课程名（页眉与书签用）
%   \newcommand{\YHBversion}{v1.2}        % 版本号（封面用）
%   \newcommand{\YHBquote}{……}            % 封面题句
%   \definecolor{yhbAccent}{HTML}{24465C} % 主色（深）
%   \definecolor{yhbAccent2}{HTML}{347A80}% 辅色（浅）
%   \YHBbooktrue / \YHBbookfalse          % 是否 book 文档类（有 \chapter）
% 章节文件只写内容，不要重复 \usepackage、不要改字号。
% ============================================================

% ---------- 页面与字体 ----------
\usepackage{geometry}
\geometry{a4paper,left=24mm,right=24mm,top=23mm,bottom=24mm,headheight=15pt,headsep=7mm,footskip=13mm}
\usepackage{microtype}
\usepackage{amsmath,amssymb,amsfonts,mathtools}
\usepackage{esint}
\usepackage{graphicx}
\usepackage{adjustbox}
\usepackage{array,booktabs,multirow,makecell,tabularx,longtable}
\usepackage{xcolor,colortbl}
\usepackage{enumitem}
\usepackage{needspace,etoolbox}
\usepackage{titlesec,titletoc}
\usepackage{fancyhdr}
\usepackage{caption}
\usepackage[most]{tcolorbox}

% ---------- 颜色别名（章节里用的旧名继续有效） ----------
\colorlet{mainblue}{yhbAccent}
\colorlet{secondcyan}{yhbAccent2}
\colorlet{thirdteal}{yhbAccent2}
\colorlet{highlight}{yhbAccent!6!white}
\colorlet{tablehead}{yhbAccent!12!white}
\definecolor{textgray}{HTML}{4A5057}

% ---------- 课程标签 ----------
\providecommand{\YHBdoclabel}{课程笔记}
% ---------- 正文节奏 ----------
\linespread{1.3}
\setlength{\parindent}{0pt}
\setlength{\parskip}{6pt plus 1.5pt minus 1.5pt}
\setlength{\emergencystretch}{2.5em}
\raggedbottom
\widowpenalty=10000
\clubpenalty=10000
\setlength{\tabcolsep}{7pt}
\renewcommand{\arraystretch}{1.3}
\allowdisplaybreaks[2]
\AtBeginDocument{%
  \setlength{\abovedisplayskip}{10pt plus 2pt minus 2pt}%
  \setlength{\belowdisplayskip}{10pt plus 2pt minus 2pt}%
  \setlength{\abovedisplayshortskip}{6pt}%
  \setlength{\belowdisplayshortskip}{6pt}%
}
\captionsetup{font=small,labelfont=bf,skip=6pt}

% ---------- 列表 ----------
\setlist{leftmargin=2em,labelsep=0.6em,itemsep=4pt,parsep=2pt,topsep=6pt,partopsep=0pt}
\setlist[itemize,2]{label=\textbullet,leftmargin=1.6em}

% ---------- 标题层级 ----------
\ifdefined\YHBbooktrue
  \titleformat{\chapter}[display]
    {\sffamily\bfseries\color{yhbAccent}}
    {\normalsize\sffamily 第\,\thechapter\,章}
    {6pt}
    {\LARGE}
    [{\vspace{5pt}{\color{yhbAccent!55}\titlerule[1pt]}}]
  \titlespacing*{\chapter}{0pt}{14pt}{24pt}
  \titleformat{\section}{\Large\sffamily\bfseries\color{yhbAccent}}{\thesection}{0.7em}{}
  \titlespacing*{\section}{0pt}{18pt plus 3pt minus 2pt}{8pt}
\else
  \titleformat{\section}{\Large\sffamily\bfseries\color{yhbAccent}}{\thesection}{0.7em}{}[{\vspace{3pt}{\color{yhbAccent!45}\titlerule[0.7pt]}}]
  \titlespacing*{\section}{0pt}{18pt plus 3pt minus 2pt}{9pt}
\fi
\titleformat{\subsection}{\large\sffamily\bfseries\color{yhbAccent!85!black}}{\thesubsection}{0.7em}{}
\titleformat{\subsubsection}{\normalsize\bfseries\color{black!85}}{\thesubsubsection}{0.7em}{}
\titlespacing*{\subsection}{0pt}{12pt plus 2pt minus 1pt}{5pt}
\titlespacing*{\subsubsection}{0pt}{9pt plus 2pt minus 1pt}{4pt}
\pretocmd{\section}{\Needspace{6\baselineskip}}{}{}
\pretocmd{\subsection}{\Needspace{5\baselineskip}}{}{}
\pretocmd{\subsubsection}{\Needspace{4\baselineskip}}{}{}

\ifdefined\YHBbooktrue
  \setcounter{secnumdepth}{2}
  \setcounter{tocdepth}{1}
\else
  \setcounter{secnumdepth}{3}
  \setcounter{tocdepth}{2}
\fi
\renewcommand{\contentsname}{目录}
% 目录编号颜色由 titletoc 的 titlecontents 控制，不再重定义 \numberline

% ---------- 目录样式 ----------
\ifdefined\YHBbooktrue
  \titlecontents{chapter}[0pt]
    {\addvspace{6pt}\sffamily\bfseries\color{yhbAccent}}
    {\contentslabel{2.2em}}{}{\titlerule*[0.6em]{.}\contentspage}
  \titlecontents{section}[2.2em]
    {\addvspace{2pt}}
    {\contentslabel{2.2em}}{}{\titlerule*[0.6em]{.}\contentspage}
\else
  \titlecontents{section}[0pt]
    {\addvspace{5pt}\sffamily\bfseries\color{yhbAccent}}
    {\contentslabel{2.2em}}{}{\titlerule*[0.6em]{.}\contentspage}
  \titlecontents{subsection}[2.2em]
    {\addvspace{2pt}}
    {\contentslabel{2.4em}}{}{\titlerule*[0.6em]{.}\contentspage}
\fi

% ---------- 页眉页脚 ----------
\pagestyle{fancy}
\fancyhf{}
\fancyhead[L]{\footnotesize\sffamily\color{textgray}\YHBcourse\ \YHBdoclabel}
\ifdefined\YHBbooktrue
  \fancyhead[R]{\footnotesize\sffamily\color{textgray}\nouppercase{\leftmark}}
  \renewcommand{\chaptermark}[1]{\markboth{#1}{}}
\else
  \fancyhead[R]{\footnotesize\sffamily\color{textgray}\nouppercase{\rightmark}}
  \renewcommand{\sectionmark}[1]{\markright{#1}}
\fi
\fancyfoot[C]{\footnotesize\color{textgray}\thepage}
\renewcommand{\headrulewidth}{0.35pt}
\renewcommand{\headrule}{{\color{yhbAccent!45}\hrule width\headwidth height\headrulewidth}}
\fancypagestyle{plain}{%
  \fancyhf{}%
  \fancyfoot[C]{\footnotesize\color{textgray}\thepage}%
  \renewcommand{\headrulewidth}{0pt}%
}

% ---------- 公式与考点框 ----------
\newcommand{\keyformula}[1]{%
  \begin{tcolorbox}[enhanced,breakable,colback=highlight,colframe=yhbAccent!30!white,boxrule=0.5pt,arc=2.5pt,left=11pt,right=11pt,top=9pt,bottom=9pt,boxsep=0pt,before skip=11pt,after skip=11pt,halign=center,width=\linewidth]
  \adjustbox{max width=\linewidth-2em}{$\displaystyle #1$}
  \end{tcolorbox}}
\newcommand{\corepoint}[1]{%
  \begin{tcolorbox}[enhanced,breakable,colback=highlight,colframe=yhbAccent2,boxrule=0pt,borderline west={2.4pt}{0pt}{yhbAccent2},arc=0pt,left=11pt,right=11pt,top=8pt,bottom=8pt,before skip=11pt,after skip=11pt]
  \textbf{\color{yhbAccent}核心考点}\quad #1
  \end{tcolorbox}}
\newcommand{\symb}[2]{\textbf{#1}：#2\par}

% ---------- 超链接（务必最后加载） ----------
\usepackage[unicode=true,colorlinks=true,linkcolor=yhbAccent,citecolor=yhbAccent,urlcolor=yhbAccent!85!black,pdfborder={0 0 0},bookmarksnumbered=true,pdfstartview=FitH]{hyperref}
\usepackage{hypcap}
\hypersetup{pdftitle={\YHBcourse 课程笔记},pdfauthor={浮萍},pdfsubject={课程学习与复习笔记},pdfcreator={XeLaTeX}}

\usepackage{multicol}
\usepackage{setspace}
\usepackage{amsthm}
\newtheoremstyle{yhbthm}{4pt}{5pt}{\normalfont}{0pt}{\bfseries}{：}{0.3em}{}
\theoremstyle{yhbthm}
\newtheorem{example}{例}[section]

% ---------- 本课程常用符号 ----------
\newcommand{\Cov}{\mathrm{Cov}}
\newcommand{\E}{\mathrm{E}}
\newcommand{\D}{\mathrm{Var}}
\newcommand{\Prob}{\mathrm{P}}
\newcommand{\dif}{\mathop{}\!\mathrm{d}}
\newcommand{\e}{\mathrm{e}}
\newcommand{\indep}{\perp \!\!\! \perp}
\newcommand{\samplemean}[1]{\overline{#1}}
\newcommand{\samplevar}[1]{S_{#1}^2}
\newcommand{\symnote}[1]{\par\noindent\textbf{符号说明：}#1\par}
\newcommand{\maxrv}[2]{\max\left\{#1,#2\right\}}
\newcommand{\minrv}[2]{\min\left\{#1,#2\right\}}
\newcommand{\ualpha}{u_{\alpha/2}}
\newcommand{\talpha}[1]{t_{\alpha/2}(#1)}
\newcommand{\xalpha}[1]{\chi^2_{\alpha/2}(#1)}
\newcommand{\xalphaone}[1]{\chi^2_{1-\alpha/2}(#1)}
\newcommand{\makeCover}{% ============================================================
% 2.cover.tex —— 统一封面（信息全部来自主文件定义）
% 需要的宏：\YHBcourse、\YHBversion、\YHBquote、yhbAccent、yhbAccent2
% 可选：\YHBcoverSize / \YHBcoverLead（课程名字号与行距）、\YHBcoversub（副标题）
% ============================================================
\providecommand{\YHBcoverSize}{38pt}
\providecommand{\YHBcoverLead}{46pt}
\providecommand{\YHBcoversub}{课程笔记}
\begin{titlepage}
\thispagestyle{empty}
\pdfbookmark[0]{封面}{cover}

\begin{tikzpicture}[remember picture,overlay]
  % 顶部深色块
  \fill[yhbAccent] (current page.north west) rectangle ([yshift=-132mm]current page.north east);
  % 分隔细条（辅色）
  \fill[yhbAccent2] ([yshift=-132mm]current page.north west) rectangle ([yshift=-135mm]current page.north east);
  % 极细横线，弱装饰，不压字
  \foreach \y in {20,34,48,62,76,90,104,118}{%
    \draw[white,opacity=0.08,line width=0.5pt]
      ([xshift=16mm,yshift=-\y mm]current page.north west) --
      ([xshift=-16mm,yshift=-\y mm]current page.north east);}
\end{tikzpicture}

\vspace*{32mm}
\noindent
{\xeCJKsetup{CJKecglue={}}\color{white}\sffamily\bfseries\fontsize{\YHBcoverSize}{\YHBcoverLead}\selectfont \YHBcourse\par}
\vspace{6mm}
{\color{white}\sffamily\fontsize{25}{32}\selectfont \YHBcoversub\par}
\vspace{9mm}
{\color{white}\fontsize{12.5}{18}\selectfont 作者\quad 浮萍\qquad 版本\quad \YHBversion\par}

\vspace{20mm}
{\sffamily\bfseries\fontsize{13}{20}\selectfont\color{yhbAccent}\YHBquote\par}

\vfill
{\color{yhbAccent}\rule{\textwidth}{1.4pt}}\par
\vspace{7mm}
{\fontsize{11.5}{17}\selectfont\sffamily
\textbf{QQ 邮箱}\quad\href{mailto:3071974740@qq.com}{3071974740@qq.com}\par
\vspace{2mm}
\textbf{GitHub}\quad\href{https://github.com/Duckweed-yhb}{github.com/Duckweed-yhb}\par}
\vspace{7mm}
{\fontsize{10.5}{16}\selectfont\color{black!78}
本笔记为个人学习整理，内容仍在持续完善。如发现内容、公式、文字或排版问题，欢迎通过上述 QQ 邮箱反馈；请尽量注明版本、页码与问题描述，便于核对和修订。\par}
\vspace*{12mm}
\end{titlepage}
}

\begin{document}

% 1. 封面（无页码）
\makeCover
\thispagestyle{empty}
\newpage

% 2. 目录页（罗马数字页码，纯页码样式）
\pagestyle{plain}
\pagenumbering{Roman}
\tableofcontents
\newpage

% 3. 正文部分（阿拉伯数字页码，页眉页脚生效）
\pagestyle{fancy}
\pagenumbering{arabic}
\setcounter{page}{1}

% ===================== 第1章 随机事件与概率 =====================
\chapter{随机事件与概率}
\section{条件概率与独立性}

\begin{example}[条件概率真假命题辨析]
下列命题中，假命题是（）。

A. 若 $P(A|B) = P(A)$，则 $P(A|\overline{B}) = P(A)$

B. 若 $P(A|B) > P(A)$，则 $P(\overline{A}|\overline{B}) > P(\overline{A})$

C. 若 $P(A|B) > P(A|\overline{B})$，则 $P(A) > P(AB)$

D. 若 $P(A|A \cup B) > P(\overline{A}|A \cup B)$，则 $P(A) > P(B)$
\end{example}

\subsubsection{解}
\begin{itemize}
  \item 选项A：由 $P(A|B) = P(A)$ 知 $A,B$ 相互独立，故 $A$ 与 $\overline{B}$ 独立，$P(A|\overline{B}) = P(A)$，真命题。
  \item 选项B：由 $P(A|B) > P(A)$ 得 $P(AB) > P(A)P(B)$，结合 $P(A) = P(AB) + P(A\overline{B})$，得 $P(A\overline{B}) < P(A)P(\overline{B})$，则 $P(\overline{A}|\overline{B}) = 1 - P(A|\overline{B}) > 1 - P(A) = P(\overline{A})$，真命题。
  \item 选项C：取反例 $A \subset B$，则 $P(A|B) = \frac{P(A)}{P(B)} > 0 = P(A|\overline{B})$（$P(B) < 1$），且 $AB = A$，故 $P(A) = P(AB)$，不满足 $P(A) > P(AB)$，假命题。
  \item 选项D：$P(A|A \cup B) = \frac{P(A)}{P(A \cup B)}$，$P(\overline{A}|A \cup B) = \frac{P(B - A)}{P(A \cup B)}$，由不等式得 $P(A) > P(B - A)$；又 $P(B) = P(AB) + P(B - A) \le P(A) + P(B - A)$，故 $P(A) > P(B)$，真命题。
\end{itemize}

最终答案：$\boxed{C}$

\begin{example}[独立事件与条件概率的综合应用]
设 $A$，$B$ 事件满足 $P(AB) > P(A)P(B)$，若存在事件 $C \subset AB$，使得 $A-C$ 与 $B$ 独立，那么 $P(C)=$（）。

A. $P(A) - P(A|\overline{B})$ \quad\quad B. $P(A) - P(A|B)$

C. $P(B) - P(B|\overline{A})$ \quad\quad D. $P(B) - P(B|A)$
\end{example}

\subsubsection{解}
已知 $C \subset AB$，故 $A-C = A \cap \overline{C}$，且 $(A-C) \cap B = AB - C$。
由 $A-C$ 与 $B$ 独立，得
\[
P((A-C) \cap B) = P(A-C)P(B).
\]
代入 $P((A-C) \cap B) = P(AB) - P(C)$ 和 $P(A-C) = P(A) - P(C)$，得
\[
P(AB) - P(C) = \bigl(P(A) - P(C)\bigr)P(B).
\]
展开并整理：
\[
P(AB) - P(C) = P(A)P(B) - P(C)P(B),
\]
\[
P(AB) - P(A)P(B) = P(C)\bigl(1 - P(B)\bigr),
\]
\[
P(C) = \frac{P(AB) - P(A)P(B)}{1 - P(B)}.
\]

对选项逐一验证：
\begin{itemize}
  \item 选项A：$P(A) - P(A|\overline{B}) = P(A) - \frac{P(A\overline{B})}{P(\overline{B})} = P(A) - \frac{P(A) - P(AB)}{1 - P(B)}$
  \[
  = \frac{P(A)(1 - P(B)) - P(A) + P(AB)}{1 - P(B)} = \frac{P(AB) - P(A)P(B)}{1 - P(B)},
  \]
  与 $P(C)$ 完全一致。
  \item 选项B：$P(A) - P(A|B) = P(A) - \frac{P(AB)}{P(B)}$，与 $P(C)$ 形式不同。
  \item 选项C：$P(B) - P(B|\overline{A}) = \frac{P(AB) - P(A)P(B)}{1 - P(A)}$，分母为 $1 - P(A)$，与 $P(C)$ 不同。
  \item 选项D：$P(B) - P(B|A) = P(B) - \frac{P(AB)}{P(A)}$，与 $P(C)$ 形式不同。
\end{itemize}

最终答案：$\boxed{A}$

% ===================== 新增例题：分布函数积分证明 =====================
\begin{example}[分布函数积分恒等式证明]
设 $F(x)$ 是连续型随机变量 $X$ 的分布函数，且存在二阶导数，常数 $a>0$。证明：
\[
\int_{-\infty}^{+\infty} \bigl[F(x+a) - F(x)\bigr] \mathrm{d}x = a
\]
\end{example}

\subsubsection{证明}
记 $X$ 的概率密度为 $f(x) = F'(x)$。

1.  \textbf{改写被积函数}：
    由分布函数定义，$F(x+a) - F(x) = P(x < X \le x+a)$。
    因此，被积函数可表示为概率密度的积分：
    \[
    F(x+a) - F(x) = \int_{x}^{x+a} f(t) \mathrm{d}t.
    \]

2.  \textbf{交换积分次序}：
    将上述表达式代入原积分，并交换积分次序（由Fubini定理保证合法）：
    \[
    \begin{aligned}
    \int_{-\infty}^{+\infty} \bigl[F(x+a) - F(x)\bigr] \mathrm{d}x
    &= \int_{-\infty}^{+\infty} \left( \int_{x}^{x+a} f(t) \mathrm{d}t \right) \mathrm{d}x \\
    &= \int_{-\infty}^{+\infty} f(t) \left( \int_{t-a}^{t} \mathrm{d}x \right) \mathrm{d}t.
    \end{aligned}
    \]

3.  \textbf{计算内层积分}：
    对 $x$ 积分：
    \[
    \int_{t-a}^{t} \mathrm{d}x = (t) - (t-a) = a.
    \]

4.  \textbf{完成证明}：
    将内层积分结果代入，得到：
    \[
    \int_{-\infty}^{+\infty} \bigl[F(x+a) - F(x)\bigr] \mathrm{d}x
    = \int_{-\infty}^{+\infty} f(t) \cdot a \, \mathrm{d}t
    = a \int_{-\infty}^{+\infty} f(t) \mathrm{d}t.
    \]
    由于 $f(t)$ 是概率密度，满足 $\int_{-\infty}^{+\infty} f(t) \mathrm{d}t = 1$，因此：
    \[
    \int_{-\infty}^{+\infty} \bigl[F(x+a) - F(x)\bigr] \mathrm{d}x = a \cdot 1 = a.
    \]
    证毕。

% ===================== 新增例题结束 =====================

% ===================== 第2章 多维随机变量及其分布 =====================
\chapter{多维随机变量及其分布}
\section{随机变量函数的分布}

\begin{example}[离散型随机变量函数的分布函数]
已知随机变量 $X$ 的分布函数为
\[
F(x) =
\begin{cases}
0, & x < -1, \\
0.2, & -1 \le x < 0, \\
0.6, & 0 \le x < 1, \\
1, & x \ge 1,
\end{cases}
\]
令 $Y = X^2$，求 $Y$ 的分布函数 $F_Y(y)$。
\end{example}

\subsubsection{解}
首先，由分布函数 $F(x)$ 的跳跃点可知，$X$ 为离散型随机变量，其可能取值为 $-1, 0, 1$，对应概率为分布函数的跳跃高度：
\[
\begin{aligned}
P(X = -1) &= F(-1) - F(-1^-) = 0.2 - 0 = 0.2, \\
P(X = 0) &= F(0) - F(0^-) = 0.6 - 0.2 = 0.4, \\
P(X = 1) &= F(1) - F(1^-) = 1 - 0.6 = 0.4.
\end{aligned}
\]
其次，由 $Y = X^2$ 可知，$Y$ 的可能取值为 $0$ 和 $1$，对应概率为：
\[
\begin{aligned}
P(Y = 0) &= P(X = 0) = 0.4, \\
P(Y = 1) &= P(X = -1) + P(X = 1) = 0.2 + 0.4 = 0.6.
\end{aligned}
\]
最后，根据分布函数的定义 $F_Y(y) = P(Y \le y)$，分区间讨论：
\begin{itemize}
  \item 当 $y < 0$ 时，$Y = X^2 \ge 0$，故 $P(Y \le y) = 0$，即 $F_Y(y) = 0$。
  \item 当 $0 \le y < 1$ 时，$Y \le y$ 等价于 $Y = 0$，故 $P(Y \le y) = P(Y = 0) = 0.4$，即 $F_Y(y) = 0.4$。
  \item 当 $y \ge 1$ 时，$Y \le y$ 包含 $Y = 0$ 和 $Y = 1$，故 $P(Y \le y) = P(Y = 0) + P(Y = 1) = 1$，即 $F_Y(y) = 1$。
\end{itemize}
因此，$Y$ 的分布函数为
\[
\boldsymbol{F_Y(y)} =
\begin{cases}
0, & y < 0, \\
0.4, & 0 \le y < 1, \\
1, & y \ge 1.
\end{cases}
\]

\begin{example}[混合分布函数的概率密度]
设随机变量 $X,Y$ 相互独立，且 $X$ 服从标准正态分布 $N(0,1)$，$Y$ 的概率分布为 $P(Y=1) = P(Y=2) = \frac{1}{2}$，求 $Z = X^Y$ 的概率密度函数 $f_Z(z)$。
\end{example}

\subsubsection{解}
我们使用全概率公式来求解 $Z = X^Y$ 的分布函数 $F_Z(z) = P(Z \le z)$。
由全概率公式及 $X,Y$ 相互独立，可得：
\[
\begin{aligned}
F_Z(z) &= P(Z \le z) = P(X^Y \le z) \\
&= P(Y=1)P(X^1 \le z \mid Y=1) + P(Y=2)P(X^2 \le z \mid Y=2) \\
&= \frac{1}{2}P(X \le z) + \frac{1}{2}P(X^2 \le z).
\end{aligned}
\]
记标准正态分布的分布函数为 $\Phi(z)$，概率密度为 $\varphi(z) = \frac{1}{\sqrt{2\pi}}e^{-z^2/2}$。分情况讨论：
\begin{itemize}
  \item 当 $z < 0$ 时：
    $P(X \le z) = \Phi(z)$，而 $P(X^2 \le z) = 0$，故
    \[
    F_Z(z) = \frac{1}{2}\Phi(z).
    \]
  \item 当 $z \ge 0$ 时：
    $P(X \le z) = \Phi(z)$，且 $P(X^2 \le z) = P(-\sqrt{z} \le X \le \sqrt{z}) = \Phi(\sqrt{z}) - \Phi(-\sqrt{z}) = 2\Phi(\sqrt{z}) - 1$，故
    \[
    F_Z(z) = \frac{1}{2}\Phi(z) + \frac{1}{2}\left(2\Phi(\sqrt{z}) - 1\right) = \frac{1}{2}\Phi(z) + \Phi(\sqrt{z}) - \frac{1}{2}.
    \]
\end{itemize}

对分布函数 $F_Z(z)$ 求导，即可得到概率密度函数 $f_Z(z)$：
\begin{itemize}
  \item 当 $z < 0$ 时：
    \[
    f_Z(z) = \frac{1}{2}\varphi(z) = \frac{1}{2\sqrt{2\pi}}e^{-z^2/2}.
    \]
  \item 当 $z > 0$ 时：
    \[
    \begin{aligned}
    f_Z(z) &= \frac{1}{2}\varphi(z) + \varphi(\sqrt{z}) \cdot \frac{1}{2\sqrt{z}} \\
    &= \frac{1}{2\sqrt{2\pi}}e^{-z^2/2} + \frac{1}{\sqrt{2\pi}}e^{-z/2} \cdot \frac{1}{2\sqrt{z}} \\
    &= \frac{1}{2\sqrt{2\pi}}e^{-z^2/2} + \frac{1}{2\sqrt{2\pi z}}e^{-z/2}.
    \end{aligned}
    \]
  \item 当 $z = 0$ 时，密度可定义为 $0$ 或不单独定义（单点不影响积分）。
\end{itemize}

因此，$Z$ 的概率密度函数为
\[
\boldsymbol{f_Z(z)} =
\begin{cases}
\displaystyle \frac{1}{2\sqrt{2\pi}}e^{-z^2/2}, & z < 0, \\[6pt]
\displaystyle \frac{1}{2\sqrt{2\pi}}e^{-z^2/2} + \frac{1}{2\sqrt{2\pi z}}e^{-z/2}, & z > 0, \\[6pt]
0, & z = 0.
\end{cases}
\]

% 新增例题：泊松+指数分布期望计算
\begin{example}[泊松分布与指数分布的期望综合计算]
设随机变量 $X \sim P(\lambda)$（泊松分布），$Y \sim E(\lambda)$（指数分布），且 $X$ 与 $Y$ 相互独立。已知 $E(X) = E(Y)$，求：\\
(1) 参数 $\lambda$ 的值；\\
(2) 期望 $E(X^2 2^Y)$。
\end{example}

\subsubsection{解}
(1) 求参数 $\lambda$：
由分布的期望性质：
- $X \sim P(\lambda)$，故 $E(X) = \lambda$；
- $Y \sim E(\lambda)$，故 $E(Y) = \frac{1}{\lambda}$。

根据题意 $E(X) = E(Y)$，得方程：
\[
\lambda = \frac{1}{\lambda} \implies \lambda^2 = 1.
\]
因分布参数 $\lambda > 0$，故 $\lambda = 1$。

(2) 计算 $E(X^2 2^Y)$：
由 $X$ 与 $Y$ 相互独立，得 $E(X^2 2^Y) = E(X^2) \cdot E(2^Y)$，分别计算两部分：

(1) 计算 $E(X^2)$：
当 $\lambda = 1$ 时，$X \sim P(1)$，其期望 $E(X) = 1$，方差 $D(X) = 1$。
由方差公式 $D(X) = E(X^2) - [E(X)]^2$，得：
\[
E(X^2) = D(X) + [E(X)]^2 = 1 + 1^2 = 2.
\]

(2) 计算 $E(2^Y)$：
当 $\lambda = 1$ 时，$Y \sim E(1)$，其概率密度为 $f_Y(y) = e^{-y},\ y > 0$。
由连续型随机变量函数的期望公式：
\[
\begin{aligned}
E(2^Y) &= \int_{0}^{+\infty} 2^y e^{-y} \,\mathrm{d}y \\
&= \int_{0}^{+\infty} e^{y \ln 2} e^{-y} \,\mathrm{d}y \\
&= \int_{0}^{+\infty} e^{-y(1 - \ln 2)} \,\mathrm{d}y \\
&= \frac{1}{1 - \ln 2}.
\end{aligned}
\]

(3) 合并结果：
\[
E(X^2 2^Y) = E(X^2) \cdot E(2^Y) = 2 \cdot \frac{1}{1 - \ln 2} = \frac{2}{1 - \ln 2}.
\]

最终结果：
(1) $\boxed{\lambda = 1}$；
(2) $\boxed{E(X^2 2^Y) = \dfrac{2}{1 - \ln 2}}$。

\section{边缘分布}

\begin{example}[二维连续型随机变量的边缘分布]
已知二维随机变量 $(X,Y)$ 的联合概率密度为
\[
f(x,y) = 
\begin{cases}
x \mathrm{e}^{-x(1+y)}, & x>0,\ y>0, \\
0, & \text{其他}.
\end{cases}
\]
分别求 $X$ 与 $Y$ 的边缘分布概率密度。
\end{example}

\subsubsection{解}
1. 求 $f_X(x)$：
当 $x > 0$ 时，
\[
f_X(x) = \int_{0}^{+\infty} x \mathrm{e}^{-x(1+y)} \mathrm{d}y = x\mathrm{e}^{-x} \int_{0}^{+\infty} \mathrm{e}^{-xy} \mathrm{d}y = x\mathrm{e}^{-x} \cdot \frac{1}{x} = \mathrm{e}^{-x}.
\]
当 $x \le 0$ 时，$f_X(x) = 0$，故
\[
f_X(x) = \begin{cases} \mathrm{e}^{-x}, & x>0, \\ 0, & x\le0. \end{cases}
\]

2. 求 $f_Y(y)$：
当 $y > 0$ 时，
\[
f_Y(y) = \int_{0}^{+\infty} x \mathrm{e}^{-x(1+y)} \mathrm{d}x = \frac{1}{(1+y)^2}.
\]
当 $y \le 0$ 时，$f_Y(y) = 0$，故
\[
f_Y(y) = \begin{cases} \displaystyle \frac{1}{(1+y)^2}, & y>0, \\ 0, & x\le0. \end{cases}
\]

\begin{example}[均匀分布条件下的联合/边缘分布与概率计算]
从区间 $(1,3)$ 中随机取一个数，记为 $X$，再从 $(X,3)$ 中取一个数，设为 $Y$。求：\\
(1) 联合概率密度 $f(x,y)$；\\
(2) $Y$ 的边缘概率密度 $f_Y(y)$；\\
(3) 概率 $P(X + Y \leq 4)$。
\end{example}

\subsubsection{解}
由题意，$X \sim U(1,3)$，故 $f_X(x) = \begin{cases} \frac{1}{2}, & 1<x<3, \\ 0, & \text{其他} \end{cases}$；
条件分布 $Y|X=x \sim U(x,3)$，故 $f_{Y|X}(y|x) = \begin{cases} \frac{1}{3-x}, & x<y<3, \\ 0, & \text{其他} \end{cases}$。

(1) 联合概率密度：
\[
f(x,y) = f_X(x) \cdot f_{Y|X}(y|x) = \begin{cases} \frac{1}{2(3-x)}, & 1<x<y<3, \\ 0, & \text{其他}. \end{cases}
\]

(2) 边缘概率密度 $f_Y(y)$：
当 $1<y<3$ 时，
\[
f_Y(y) = \int_{1}^{y} \frac{1}{2(3-x)} \mathrm{d}x = \frac{1}{2} \left. \ln(3-x) \right|_{y}^{1} = \frac{1}{2} \ln \frac{2}{3-y}.
\]
当 $y \notin (1,3)$ 时，$f_Y(y) = 0$，故
\[
f_Y(y) = \begin{cases} \frac{1}{2} \ln \frac{2}{3-y}, & 1<y<3, \\ 0, & \text{其他}. \end{cases}
\]

(3) 概率 $P(X + Y \leq 4)$：
积分区域为 $1<x\le2$，$x<y\le4-x$，则
\[
P(X + Y \leq 4) = \int_{1}^{2} \int_{x}^{4-x} \frac{1}{2(3-x)} \mathrm{d}y \mathrm{d}x = \int_{1}^{2} \frac{2-x}{3-x} \mathrm{d}x.
\]
% 修复：手动换行长公式，避免Overfull \hbox
令被积函数变形为 $1 - \frac{1}{3-x}$，积分得
\[
\int_{1}^{2} \left(1 - \frac{1}{3-x}\right) \mathrm{d}x = \left. \left(x + \ln(3-x)\right) \right|_{1}^{2} = (2 + 0) - (1 + \ln2) = 1 - \ln2.
\]
最终结果：$\boxed{1 - \ln2}$。

% ===================== 新增例题：P(X+Y < 1) 计算 =====================
\begin{example}[二维连续型随机变量概率计算（X+Y < 1）]
已知二维随机变量 $(X,Y)$ 的联合概率密度为 $f(x,y) = x e^{-y}$（$0 < x < y$），求概率 $P(X+Y < 1)$。
\end{example}

\subsubsection{解}
要计算 $P(X+Y < 1)$，需在满足 $0 < x < y$ 且 $x + y < 1$ 的区域上对联合密度进行二重积分。

步骤1：确定积分区域
联立不等式 $0 < x < y$ 和 $x + y < 1$，可得积分区域：
\[
0 < x < \frac{1}{2}, \quad x < y < 1 - x
\]

步骤2：建立二重积分
\[
\begin{aligned}
P(X+Y < 1) &= \iint_{x+y<1} f(x,y) \, dxdy \\
&= \int_{0}^{\frac{1}{2}} \int_{x}^{1-x} x e^{-y} \, dy dx
\end{aligned}
\]

步骤3：计算内层积分（对 $y$ 积分）
\[
\begin{aligned}
\int_{x}^{1-x} x e^{-y} \, dy &= x \left[ -e^{-y} \right]_{x}^{1-x} \\
&= x \left( -e^{-(1-x)} + e^{-x} \right) \\
&= x e^{-x} - x e^{x-1}
\end{aligned}
\]

步骤4：计算外层积分（对 $x$ 积分）
\[
P(X+Y < 1) = \int_{0}^{\frac{1}{2}} \left( x e^{-x} - x e^{x-1} \right) dx
\]
使用分部积分法计算：
- $\int x e^{-x} dx = -x e^{-x} - e^{-x} + C$
- $\int x e^{x} dx = x e^{x} - e^{x} + C$

代入上下限 $0$ 和 $\frac{1}{2}$ 并化简：
\[
\begin{aligned}
\int_{0}^{\frac{1}{2}} x e^{-x} dx &= 1 - \frac{3}{2} e^{-\frac{1}{2}}, \\
\int_{0}^{\frac{1}{2}} x e^{x} dx &= 1 - \frac{1}{2} e^{\frac{1}{2}}.
\end{aligned}
\]

合并结果：
\[
\begin{aligned}
P(X+Y < 1) &= \left( 1 - \frac{3}{2} e^{-\frac{1}{2}} \right) - e^{-1} \left( 1 - \frac{1}{2} e^{\frac{1}{2}} \right) \\
&= 1 - e^{-1} - e^{-\frac{1}{2}}.
\end{aligned}
\]

最终答案：$\boxed{P(X+Y < 1) = 1 - e^{-1} - e^{-\frac{1}{2}}}$

% ===================== 新增例题结束 =====================

\begin{example}[均匀分布下最大值与最小值的期望计算]
在 $[0,1]$ 内任取两点 $X,Y$，记 $U = \max\{X,Y\}$，$V = \min\{X,Y\}$，求 $E[(1-U)(1-V)]$。
\end{example}

\subsubsection{解}
由题意，$X,Y$ 独立同分布于 $U(0,1)$，故联合密度为
\[
f(x,y) = \begin{cases}
1, & 0 \le x,y \le 1, \\
0, & \text{其他}.
\end{cases}
\]
注意到 $(1-U)(1-V) = (1-\max\{X,Y\})(1-\min\{X,Y\}) = (1-X)(1-Y)$，因此
\[
E[(1-U)(1-V)] = E[(1-X)(1-Y)] = E(1 - X - Y + XY).
\]
由独立性，$E(XY) = E(X)E(Y)$，且 $E(X) = E(Y) = \frac{1}{2}$，代入得
\[
E[(1-U)(1-V)] = 1 - E(X) - E(Y) + E(X)E(Y) = 1 - \frac{1}{2} - \frac{1}{2} + \frac{1}{2} \cdot \frac{1}{2} = \frac{1}{4}.
\]
最终结果：$\boxed{\dfrac{1}{4}}$。

% ===================== 第5章 数理统计的基本概念 =====================
\chapter{数理统计的基本概念}
\section{样本均值与方差的分布}

\begin{example}[正态总体样本均值与方差的方差计算]
设 $X_1,X_2,\dots,X_n$ 是取自总体 $X \sim N(0,\sigma^2)$ 的样本，$\overline{X}$ 和 $S^2$ 分别为样本均值和样本方差，求 $D\left[\overline{X}^2 + \left(1 - \frac{1}{n}\right)S^2\right]$。
\end{example}

\subsubsection{解}
核心性质：$\overline{X}$ 与 $S^2$ 相互独立，故
\[
D\left[\overline{X}^2 + \left(1 - \frac{1}{n}\right)S^2\right] = D(\overline{X}^2) + \left(1 - \frac{1}{n}\right)^2 D(S^2).
\]

步骤1：计算 $D(\overline{X}^2)$
$\overline{X} \sim N\left(0, \frac{\sigma^2}{n}\right)$，令 $Z = \frac{\overline{X}}{\sigma/\sqrt{n}} \sim N(0,1)$，则 $\overline{X}^2 = \frac{\sigma^2}{n} Z^2$。
因 $Z^2 \sim \chi^2(1)$，故 $D(Z^2) = 2$，因此
\[
D(\overline{X}^2) = \left(\frac{\sigma^2}{n}\right)^2 \cdot D(Z^2) = \frac{2\sigma^4}{n^2}.
\]

步骤2：计算 $D(S^2)$
由 $\frac{(n-1)S^2}{\sigma^2} \sim \chi^2(n-1)$，记 $W = \frac{(n-1)S^2}{\sigma^2}$，则 $S^2 = \frac{\sigma^2}{n-1} W$。
因 $D(W) = 2(n-1)$，故
\[
D(S^2) = \left(\frac{\sigma^2}{n-1}\right)^2 \cdot D(W) = \frac{\sigma^4}{(n-1)^2} \cdot 2(n-1) = \frac{2\sigma^4}{n-1}.
\]

步骤3：合并结果
\[
D\left[\overline{X}^2 + \left(1 - \frac{1}{n}\right)S^2\right] = \frac{2\sigma^4}{n^2} + \left(\frac{n-1}{n}\right)^2 \cdot \frac{2\sigma^4}{n-1}.
\]
% 修复：手动换行长公式，避免Overfull \hbox
化简得
\[
\frac{2\sigma^4}{n^2} + \frac{2(n-1)\sigma^4}{n^2} = \frac{2n\sigma^4}{n^2} = \frac{2\sigma^4}{n}.
\]
最终结果：$\boxed{\dfrac{2\sigma^4}{n}}$。

\begin{example}[泊松分布样本均值的概率计算]
已知随机变量 $X$ 服从参数为 $1$ 的泊松分布 $P(1)$，$X_1,X_2,X_3,X_4$ 为取自 $X$ 的简单随机样本，则 $P\left(\overline{X} > \frac{1}{4}\right) = \ ?$
\end{example}

\subsubsection{解}
样本均值 $\overline{X} = \frac{1}{4}(X_1+X_2+X_3+X_4)$，故
\[
\overline{X} > \frac{1}{4} \iff X_1+X_2+X_3+X_4 > 1.
\]
由泊松分布的可加性，$X_1+X_2+X_3+X_4 \sim P(4)$，记 $Y = X_1+X_2+X_3+X_4$，则
\[
P\left(\overline{X} > \frac{1}{4}\right) = P(Y > 1) = 1 - P(Y=0) - P(Y=1).
\]
泊松分布概率公式：
\[
P(Y=k) = \frac{4^k e^{-4}}{k!}, \quad k=0,1,2,\dots
\]
代入 $k=0,1$：
\[
P(Y=0) = e^{-4}, \quad P(Y=1) = 4e^{-4}.
\]
因此
\[
P\left(\overline{X} > \frac{1}{4}\right) = 1 - e^{-4} - 4e^{-4} = 1 - 5e^{-4}.
\]
最终结果：$\boxed{1 - 5e^{-4}}$。

% ===================== 第6章 参数估计 =====================
\chapter{参数估计}
\section{点估计}

\begin{example}[正态总体下的概率密度、极大似然估计与无偏性]
已知 $X \sim N(0,\sigma^2)$，$Y \sim N(0,2\sigma^2)$，且 $X,Y$ 相互独立，设 $Z = X - Y$，$Z_1,Z_2,\dots,Z_n$ 是取自总体 $Z$ 的一个样本，求：\\
(1) $Z$ 的概率密度；\\
(2) $\sigma^2$ 的极大似然估计 $\hat{\sigma^2}$；\\
(3) 判断 $\hat{\sigma^2}$ 是否是 $\sigma^2$ 的无偏估计。
\end{example}

\subsubsection{解}
(1) 求 $Z$ 的概率密度：
由正态分布的线性性质，$Z = X - Y$ 仍服从正态分布。
期望：$E(Z) = E(X) - E(Y) = 0 - 0 = 0$。
方差：$D(Z) = D(X) + D(Y) = \sigma^2 + 2\sigma^2 = 3\sigma^2$。
因此 $Z \sim N(0, 3\sigma^2)$，其概率密度为
\[
f_Z(z) = \frac{1}{\sqrt{2\pi \cdot 3\sigma^2}} \exp\left\{-\frac{z^2}{2 \cdot 3\sigma^2}\right\} = \frac{1}{\sqrt{6\pi}\sigma} \exp\left\{-\frac{z^2}{6\sigma^2}\right\},\quad z \in \mathbb{R}.
\]

(2) 求 $\sigma^2$ 的极大似然估计 $\hat{\sigma^2}$：
似然函数：
\[
L(\sigma^2) = \prod_{i=1}^n f_Z(Z_i) = \prod_{i=1}^n \frac{1}{\sqrt{6\pi}\sigma} \exp\left\{-\frac{Z_i^2}{6\sigma^2}\right\}
= (6\pi\sigma^2)^{-n/2} \exp\left\{-\frac{1}{6\sigma^2} \sum_{i=1}^n Z_i^2\right\}.
\]
取对数：
\[
\ln L(\sigma^2) = -\frac{n}{2}\ln(6\pi) - \frac{n}{2}\ln(\sigma^2) - \frac{1}{6\sigma^2} \sum_{i=1}^n Z_i^2.
\]
对 $\sigma^2$ 求导并令导数为0：
\[
\frac{\mathrm{d}\ln L}{\mathrm{d}(\sigma^2)} = -\frac{n}{2\sigma^2} + \frac{1}{6(\sigma^2)^2} \sum_{i=1}^n Z_i^2 = 0.
\]
解得
\[
\hat{\sigma^2} = \frac{1}{3n} \sum_{i=1}^n Z_i^2.
\]

(3) 判断无偏性：
计算期望 $E(\hat{\sigma^2})$：
\[
E(\hat{\sigma^2}) = \frac{1}{3n} \sum_{i=1}^n E(Z_i^2) = \frac{1}{3n} \cdot n \cdot E(Z^2).
\]
因 $Z \sim N(0, 3\sigma^2)$，故 $E(Z^2) = D(Z) + [E(Z)]^2 = 3\sigma^2$。
代入得
\[
E(\hat{\sigma^2}) = \frac{1}{3n} \cdot n \cdot 3\sigma^2 = \sigma^2.
\]
因此 $\hat{\sigma^2}$ 是 $\sigma^2$ 的无偏估计。

最终结果：
(1) $\boxed{f_Z(z) = \dfrac{1}{\sqrt{6\pi}\sigma} \exp\left\{-\dfrac{z^2}{6\sigma^2}\right\},\ z \in \mathbb{R}}$；
(2) $\boxed{\hat{\sigma^2} = \dfrac{1}{3n} \sum_{i=1}^n Z_i^2}$；
(3) $\hat{\sigma^2}$ \textbf{是} $\sigma^2$ 的无偏估计。

\begin{example}[最大似然估计与无偏性]
设总体 $X$ 的概率密度为
\[
f(x;\theta) = 
\begin{cases}
\displaystyle \frac{1}{\theta^2} x e^{-x/\theta}, & x>0, \\
0, & \text{其他},
\end{cases}
\]
其中 $\theta > 0$ 为未知参数。$X_1,X_2,\dots,X_n$ 是来自总体 $X$ 的简单随机样本。

1. 求 $\theta$ 的最大似然估计 $\hat\theta$；
2. 判断 $\hat\theta$ 是否为 $\theta$ 的无偏估计。
\end{example}

\subsubsection{解}
1. 最大似然估计：
似然函数：
\[
L(\theta) = \prod_{i=1}^n f(x_i;\theta) = \frac{1}{\theta^{2n}} \left(\prod_{i=1}^n x_i\right) e^{-\frac{1}{\theta}\sum_{i=1}^n x_i}.
\]
取对数：
\[
\ln L(\theta) = -2n\ln\theta + \sum_{i=1}^n \ln x_i - \frac{1}{\theta} \sum_{i=1}^n x_i.
\]
求导并令导数为0：
\[
\frac{\mathrm{d}\ln L}{\mathrm{d}\theta} = -\frac{2n}{\theta} + \frac{1}{\theta^2} \sum_{i=1}^n x_i = 0 \implies \hat\theta = \frac{1}{2n} \sum_{i=1}^n X_i = \frac{\overline{X}}{2}.
\]

2. 无偏性判断：
计算总体期望：
\[
E(X) = \int_{0}^{+\infty} x \cdot \frac{1}{\theta^2} x e^{-x/\theta} \mathrm{d}x = \frac{1}{\theta^2} \int_{0}^{+\infty} x^2 e^{-x/\theta} \mathrm{d}x = 2\theta.
\]
样本均值期望 $E(\overline{X}) = E(X) = 2\theta$，故
\[
E(\hat\theta) = E\left(\frac{\overline{X}}{2}\right) = \frac{1}{2} \cdot 2\theta = \theta.
\]
因此 $\hat\theta$ 是 $\theta$ 的无偏估计。

\subsection{估计量的评价标准}

\begin{example}[无偏估计参数求解]
设总体 $X \sim P(\lambda)$（泊松分布），$X_1,X_2,\dots,X_n$ 是来自总体 $X$ 的简单随机样本，若 $\frac{1}{n}\sum_{i=1}^n a^{X_i}$ 是 $\mathrm{e}^{\lambda}$ 的无偏估计，求常数 $a$ 的值。
\end{example}

\subsubsection{解}
1. 计算 $E[a^X]$：
由泊松分布的期望公式，
\[
E[a^X] = \sum_{k=0}^{\infty} a^k \cdot \frac{\lambda^k e^{-\lambda}}{k!} = e^{-\lambda} \sum_{k=0}^{\infty} \frac{(a\lambda)^k}{k!} = e^{-\lambda} \cdot e^{a\lambda} = e^{\lambda(a-1)}.
\]

2. 无偏性条件求解：
由无偏估计定义，$E\left(\frac{1}{n}\sum_{i=1}^n a^{X_i}\right) = E[a^X] = e^{\lambda}$，故
\[
e^{\lambda(a-1)} = e^{\lambda} \implies a-1 = 1 \implies a = 2.
\]
最终结果：$\boxed{2}$。

% 新增：无偏估计常数求解（c和k）
\begin{example}[无偏估计常数求解（c和k）]
设 $X_1,X_2,\dots,X_n$ 是取自正态总体 $N(\mu,\sigma^2)$ 的简单随机样本，求：\\
(1) 常数 $c$，使得 $c\sum_{i=1}^{n-1} (X_{i+1}-X_i)^2$ 为 $\sigma^2$ 的无偏估计；\\
(2) 常数 $k$，使得 $\hat{\sigma} = k\sum_{i=1}^n |X_i - \overline{X}|$ 为 $\sigma$ 的无偏估计。
\end{example}

\subsubsection{解}
(1) 求常数 $c$：
已知 $X_i \sim N(\mu, \sigma^2)$ 且相互独立，先计算 $E(X_{i+1}-X_i)^2$：
\[
\begin{aligned}
E(X_{i+1}-X_i)^2 &= D(X_{i+1}-X_i) + [E(X_{i+1}-X_i)]^2 \\
&= D(X_{i+1}) + D(X_i) + (\mu - \mu)^2 \\
&= \sigma^2 + \sigma^2 = 2\sigma^2.
\end{aligned}
\]
因此：
\[
E\left(\sum_{i=1}^{n-1} (X_{i+1}-X_i)^2\right) = (n-1) \cdot 2\sigma^2 = 2(n-1)\sigma^2.
\]
由无偏性要求 $E\left(c\sum_{i=1}^{n-1} (X_{i+1}-X_i)^2\right) = \sigma^2$，得：
\[
c \cdot 2(n-1)\sigma^2 = \sigma^2 \implies c = \frac{1}{2(n-1)}.
\]

(2) 求常数 $k$：
令 $Y_i = X_i - \overline{X}$，则 $Y_i \sim N\left(0, \frac{n-1}{n}\sigma^2\right)$，利用正态分布绝对值期望公式 $E|Z| = \sqrt{\frac{2}{\pi}}\sigma$（$Z \sim N(0,\sigma^2)$），得：
\[
E|Y_i| = \sqrt{\frac{2}{\pi}} \cdot \sigma \sqrt{\frac{n-1}{n}} = \sigma \sqrt{\frac{2(n-1)}{n\pi}}.
\]
因此：
\[
E\left(\sum_{i=1}^n |X_i - \overline{X}|\right) = n \cdot \sigma \sqrt{\frac{2(n-1)}{n\pi}} = \sigma \sqrt{\frac{2n(n-1)}{\pi}}.
\]
由无偏性要求 $E(\hat{\sigma}) = \sigma$，得：
\[
k \cdot \sigma \sqrt{\frac{2n(n-1)}{\pi}} = \sigma \implies k = \sqrt{\frac{\pi}{2n(n-1)}}.
\]

最终结果：
(1) $\boxed{c = \dfrac{1}{2(n-1)}}$；
(2) $\boxed{k = \sqrt{\dfrac{\pi}{2n(n-1)}}}$。

\end{document}